% ===========================================================================
\chapter{Análise de Anterioridades: Patentes Globais}
\label{cap_analise_patentes}
% ===========================================================================

A prospecção na base internacional de patentes \textbf{WIPO Patentscope} resultou inicialmente em 542 documentos globais. Ao restringir o universo para as invenções indexadas sob o código IPC \textbf{G16H} (\textit{Tecnologia da Informação e Comunicação aplicada à Saúde}), obteve-se um conjunto altamente qualificado de 86 documentos com aderência técnica direta ao domínio do processamento de dados hospitalares e eventos adversos.

A análise aprofundada desse conjunto revelou que a vasta maioria das patentes internacionais se concentra no controle administrativo de riscos corporativos ou na prevenção de perdas financeiras em redes hospitalares privadas. Dentre as 86 anterioridades, quatro documentos destacam-se como o estado da técnica mais próximo do SIGEA-GTT, sendo minuciosamente analisados a seguir.

% ===========================================================================
\section{Patente CN118888101 -- Chongqing Medical University}
\label{sec_patente_chongqing}
% ===========================================================================

\begin{itemize}
    \item \textbf{Identificação Oficial}: Pedido de Patente de Invenção CN118888101A, publicado em 2024 pela \textit{China National Intellectual Property Administration} (CNIPA) \cite{patente_cn118888101};
    \item \textbf{Título Oficial}: \textit{Patient safety training system and patient safety training method};
    \item \textbf{Aspectos Técnicos Relevantes}: O documento descreve um sistema computacional estruturado para capacitação de equipes de enfermagem e estudantes em diretrizes de segurança clínica. O sistema organiza módulos teóricos, identifica pontos críticos de vulnerabilidade do cuidado hospitalar e aplica testes de avaliação com pontuações padronizadas;
    \item \textbf{Diferencial Substantivo do SIGEA-GTT}: A patente chinesa ancora-se na prescrição de módulos e currículos teóricos estáticos, semelhantes a plataformas tradicionais de ensino a distância (\textit{e-learning}). Em contrapartida, o SIGEA-GTT opera sobre uma abordagem empírica e dinâmica: o discente é imerso em simulações realistas de auditoria retrospectiva baseada em prontuários complexos em formato semiestruturado \texttt{JSONB}, alimentando um \textit{dashboard} epidemiológico com indicadores do IHI em tempo real, sem aprisionar o aluno em um sequenciamento rígido de videoaulas ou questionários teóricos de múltipla escolha.
\end{itemize}

% ===========================================================================
\section{Patente CN111627533 -- Fortune Software Co., Ltd.}
\label{sec_patente_fortune}
% ===========================================================================

\begin{itemize}
    \item \textbf{Identificação Oficial}: Patente de Invenção CN111627533A, concedida em 2020 \cite{patente_cn111627533};
    \item \textbf{Título Oficial}: \textit{Hospital-wide adverse event active monitoring and management system};
    \item \textbf{Aspectos Técnicos Relevantes}: Propõe uma arquitetura de monitoramento ativo em escala hospitalar contínua. Incorpora módulos de extração de dados clínicos de prontuários eletrônicos de produção, rotinas de higienização de registros e emissão automatizada de alertas de eventos sentinela para gestores de risco;
    \item \textbf{Diferencial Substantivo do SIGEA-GTT}: Trata-se de uma solução puramente corporativa, voltada ao faturamento, conformidade regulatória e mitigação de processos judiciais de responsabilidade civil em hospitais reais. O SIGEA-GTT reutiliza o rigor taxonômico de triagem clínica do IHI, mas o reorienta para um ambiente seguro de aprendizagem e \textit{debriefing} pedagógico, blindando o discente de repercussões punitivas durante a sua curva de aprendizado formativo.
\end{itemize}

% ===========================================================================
\section{Patente CN113379384 -- Huaxun High Tech Co., Ltd.}
\label{sec_patente_huaxun}
% ===========================================================================

\begin{itemize}
    \item \textbf{Identificação Oficial}: Pedido de Patente CN113379384A, publicado em 2021 \cite{patente_cn113379384};
    \item \textbf{Título Oficial}: \textit{Medical event big data platform};
    \item \textbf{Aspectos Técnicos Relevantes}: Descreve uma infraestrutura corporativa baseada em tecnologias de \textit{Big Data} para o ciclo total de vida de notificações assistenciais: recepção em lote, análise de causa raiz e geração de painéis estatísticos gerenciais para diretorias hospitalares;
    \item \textbf{Diferencial Substantivo do SIGEA-GTT}: O foco da anterioridade chinesa é a escalabilidade de dados massivos em redes assistenciais centralizadas. O SIGEA-GTT converte a notificação e o cálculo de taxas epidemiológicas (como a taxa de eventos por 1.000 pacientes-dia) em instrumentos de desenvolvimento do raciocínio crítico, permitindo a comparação lado a lado entre as respostas do estudante e o gabarito validado pelo docente titular da disciplina.
\end{itemize}

% ===========================================================================
\section{Patente US20210074398 -- MedStar Health, Inc.}
\label{sec_patente_medstar}
% ===========================================================================

\begin{itemize}
    \item \textbf{Identificação Oficial}: Pedido de Patente US20210074398A1, depositado pelo sistema hospitalar norte-americano MedStar Health perante o USPTO e publicado na WIPO em 2021 \cite{patente_us20210074398};
    \item \textbf{Título Oficial}: \textit{Evaluation of patient safety event reports from free-text descriptions};
    \item \textbf{Aspectos Técnicos Relevantes}: Aborda o processamento automático de relatórios de segurança do paciente redigidos em linguagem natural livre. Utiliza modelos computacionais de Processamento de Linguagem Natural (PLN) e Aprendizado de Máquina para classificar a severidade do dano de forma autônoma;
    \item \textbf{Diferencial Substantivo do SIGEA-GTT}: A solução da MedStar Health substitui o julgamento humano por um classificador autônomo corporativo. O SIGEA-GTT adota a filosofia inversa: a ferramenta prioriza o protagonismo do estudante na detecção determinística de gatilhos clínicos e na justificativa textual de causalidade, funcionando como andaime pedagógico cognitivo que capacita o futuro enfermeiro ou médico a interpretar pistas clínicas complexas por si próprio.
\end{itemize}

% ===========================================================================
\section{Quadro Comparativo entre as Anterioridades e o SIGEA-GTT}
\label{sec_quadro_comparativo_patentes}
% ===========================================================================

A Tabela \ref{tab_comparativo_patentes} sumariza a comparação técnica e estrutural entre as quatro patentes internacionais e o sistema SIGEA-GTT.

\begin{table}[!ht]
\centering
\caption{Quadro Comparativo entre Anterioridades Globais de Patente e o SIGEA-GTT}
\label{tab_comparativo_patentes}
\footnotesize
\begin{tabularx}{\textwidth}{p{2.5cm} p{2.8cm} p{2.4cm} p{2.8cm} X}
\toprule
\textbf{Patente} & \textbf{Titular / País} & \textbf{Finalidade} & \textbf{Abordagem} & \textbf{Diferencial do SIGEA-GTT} \\
\midrule
\textbf{CN118888101} & Chongqing Med. Univ. \par (China) & Treinamento em segurança & Módulos teóricos estáticos (\textit{e-learning}) & Auditoria prática com casos fictícios dinâmicos em JSONB e cronômetro de 20 min. \\
\addlinespace
\textbf{CN111627533} & Fortune Software \par (China) & Monitoramento hospitalar & Controle corporativo e mitigação de risco & Ambiente simulado não-punitivo voltado à formação discente e debriefing formativo. \\
\addlinespace
\textbf{CN113379384} & Huaxun High Tech \par (China) & Big Data de eventos médicos & Notificação massiva corporativa & Transformação do cálculo de taxas em instrumento didático com espelho de gabarito. \\
\addlinespace
\textbf{US20210074398} & MedStar Health \par (Estados Unidos) & Classificação via PLN de textos livres & Automação algorítmica autônoma & Foco no desenvolvimento do discernimento clínico humano do discente sob supervisão. \\
\addlinespace
\textbf{SIGEA-GTT} & \textbf{UFS / DCOMP} \par \textbf{(Brasil, 2026)} & \textbf{Ensino prático de auditoria GTT} & \textbf{Simulação clínica, GTT e Qualidade} & \textbf{Metodologia IHI, Ishikawa, GUT, 5W2H, PDCA e taxas epidemiológicas integradas.} \\
\bottomrule
\end{tabularx}
\fonte{Elaborado pelo autor (2026).}
\end{table}
